\begin{solapartat}
\punts{0,1} Treball d'extracció o funció de treball: $W_0 = hf_0 = 3,04\times 10^{-19}\ \mathrm{J}$

\punts{0,25} Perquè hi hagi emissió cal que la freqüència sigui superior a la freqüència llindar.

\punts{0,1} Pel làser blau, $\lambda_b = 460\ \mathrm{nm}$:
\[ f_b = \frac{c}{\lambda_b} = \frac{3,00\times 10^{8}}{460\times 10^{-9}} = 6,52\times 10^{14}\ \mathrm{Hz} \]
Com que $f_b > f_0$ hi haurà emissió d'electrons.

\punts{0,1} Pel làser verd, $\lambda_v = 532\ \mathrm{nm}$:
\[ f_v = \frac{c}{\lambda_v} = \frac{3,00\times 10^{8}}{532\times 10^{-9}} = 5,64\times 10^{14}\ \mathrm{Hz} \]
Com $f_v > f_0$ hi haurà emissió d'electrons.

\punts{0,1} Pel làser infraroig, $\lambda_{IR} = 1\,080\ \mathrm{nm}$:
\[ f_{IR} = \frac{c}{\lambda_{IR}} = \frac{3,00\times 10^{8}}{1080\times 10^{-9}} = 2,78\times 10^{14}\ \mathrm{Hz} \]
Com que $f_v < f_0$ no hi haurà emissió d'electrons.

\destaca{Alternativament} es podia comparar l'energia del fotons amb el treball d'extracció, per al làser blau és $4,32\times 10^{-19}\ \mathrm{J} > W_0$ i per al làser verd $3,74\times 10^{-19}\ \mathrm{J} > W_0$, per tant, per aquests dos casos hi ha emissió d'electrons. En canvi per al làser d'infraroig $1,84\times 10^{-19}\ \mathrm{J} < W_0$, per tant, no hi ha emissió d'electrons.

\punts{0,2} Per determinar l'energia dels electrons emesos cal imposar el balanç d'energia:

$E_C = hf - W_0$.

El resultat obtingut és:

\begin{center}
\resizebox{\linewidth}{!}{%
\begin{tabular}{|c|c|c|c|c|c|}
\hline
\textit{Tipus de punter làser} & \textit{Longitud d'ona (nm)} & \textit{Freqüència ($\times 10^{14}$ Hz)} & \textit{Energia fotó ($\times 10^{-19}$ J)} & \textit{$E_c$ electró ($\times 10^{-19}$ J)} & \textit{$E_c$ electró (eV)} \\
\hline
Làser blau & 460 & 6,52 & 4,32 & 1,28 & 0,799 \\
\hline
Làser verd & 532 & 5,64 & 3,74 & 0,696 & 0,434 \\
\hline
Làser infraroig & 1\,080 & 2,78 & 1,84 & - & - \\
\hline
Llindar & & 4,59 & 3,04 & 0 & 0 \\
\hline
\end{tabular}}
\end{center}

\punts{0,4} Si no s'indiquen les unitats als eixos o un títol dels eixos, cal restar 0,25 p.

Si es representen punts de la recta per sota la freqüència llindar de valor diferent a zero, cal restar 0,25 punts.

Si la representació no és una recta, cal restar 0,25 punts.

Si els eixos no estan correctament escalats, cal restar 0,1 punts.

Si es comenten diversos error i el valor a restar és superior a 0,4 p, no es farà la resta completa, simplement la qualificació de la representació gràfica serà 0 punts, els errors en aquesta part no s'acumulen a la resta del problema.

\figura[0.75]{sol_fig1.png}
\end{solapartat}

\begin{solapartat}
\punts{0,3} Càlcul de l'energia cinètica
\[ E_C = hf - W_0 = 5,30\times 10^{-19} - 3,04\times 10^{-19} = 2,26\times 10^{-19}\ \mathrm{J} \]

\punts{0,45} I la velocitat és:
\[ v = \sqrt{2\frac{E_C}{m_e}} = 7,04\times 10^{5}\ \mathrm{m/s} \]

\punts{0,5} I la longitud de De Broglie és:
\[ \lambda = \frac{h}{m_e v} = \frac{6,63\times 10^{-34}}{9,11\times 10^{-31}\times 7,04\times 10^{5}} = 1,03\times 10^{-9}\ \mathrm{m} \]
\end{solapartat}
